\documentclass{article}

% if you need to pass options to natbib, use, e.g.:
%     \PassOptionsToPackage{numbers, compress}{natbib}
% before loading neurips_2025

% ready for submission
\usepackage[preprint]{neurips_2025}

% to compile a preprint version, e.g., for submission to arXiv, add add the
% [preprint] option:
%     \usepackage[preprint]{neurips_2025}

% to compile a camera-ready version, add the [final] option, e.g.:
%     \usepackage[final]{neurips_2025}

% to avoid loading the natbib package, add option nonatbib:
%    \usepackage[nonatbib]{neurips_2025}

\usepackage[utf8]{inputenc} % allow utf-8 input
\usepackage[T1]{fontenc}    % use 8-bit T1 fonts
\usepackage{hyperref}       % hyperlinks
\usepackage{url}            % simple URL typesetting
\usepackage{booktabs}       % professional-quality tables
\usepackage{amsfonts}       % blackboard math symbols
\usepackage{nicefrac}       % compact symbols for 1/2, etc.
\usepackage{microtype}      % microtypography
\usepackage{xcolor}         % colors
\usepackage{graphicx}       % images
\usepackage{multirow}       % for tables

\title{Entity Resolution at 500M Scale: Can Fine-Tuned Dense Retrievers Beat Lexical Baselines?}

% The \author macro works with any number of authors. There are two commands
% used to separate the names and addresses of multiple authors: \And and \AND.
%
% Using \And between authors leaves it to LaTeX to determine where to break the
% lines. Using \AND forces a line break at that point. So, if LaTeX puts 3 of 4
% authors names on the first line, and the last on the second line, try using
% \AND instead of \And before the third author name.

\author{%
  Anonymous Authors \\
  Department of Computer Science\\
  University of Anonymous\\
  \texttt{hi@anonymous.edu} \\
}

\begin{document}

\maketitle

\begin{abstract}
Entity resolution across hundreds of millions of user-generated records is fundamentally a retrieval problem. Lexical systems like BM25 excel at exact matching but degrade rapidly on corrupted, missing, or swapped fields. We evaluate whether compact dense embedding models (22M to 149M parameters), fine-tuned with Matryoshka Representation Learning (MRL) and Multiple Negatives Ranking Loss (MNRL), can outperform BM25 on realistic B2B contact query distributions while remaining within production memory limits. We assess these models via an ablation over dimensionality ranges and quantizations. We demonstrate that GTE-ModernBERT (149M) achieves a $+0.7$pp overall Recall@10 gain compared to BM25, and notably improves performance on severely corrupted queries (dropping email and company) by $+4.7$pp. However, models strictly reliant on prompt prefixes for zero-shot performance completely collapse under standard MNRL fine-tuning. Concurrently, lightweight models combined with integer quantization hold accuracy and compress inference latency and index volume significantly.
\end{abstract}

\section{Introduction}

Identifying duplicate or related people profiles at massive scale ($\sim$500 million records) requires a fast, low-footprint retrieval engine. The standard industry approach uses BM25. This lexical model matches perfectly when searching for "Jonathan Smith" with a clean work email, but breaks on abbreviations ("Jon"), typos ("Smyth"), missing fields, and swapped semantic categories (e.g., matching a personal Gmail rather than a corporate domain). 

Replacing BM25 with a dense retriever is theoretically straightforward but practically constrained. A standard 768-dimensional FP32 index for 500 million records requires over 1.5TB of memory—prohibitive for real-time application. Matryoshka Representation Learning (MRL) \cite{kusner2015matryoshka} mitigates this by enabling truncation of dense embeddings without structural retraining. 

We measure the recall and latency of fine-tuned MRL dense retrievers against BM25 across six distinct real-world corruption distributions: pristine, missing first name, missing email and company, name typos, domain mismatches, and swapped attributes.
We prove that while BM25 is unbeatable on pristine data, fine-tuned dense embeddings provide significant robustness on partial records. GTE-ModernBERT (149M) outperforms the baseline overall. Conversely, models relying on instruction prefixes, like Nomic-v1.5, exhibit catastrophic forgetting when exposed to standard MNRL fine-tuning on short structured strings. Finally, we show that simple, small architectures combined with 128-dim int8 quantization maintain robust recall thresholds while driving latency down near lexical baseline speeds.

\section{Methodology}

\subsection{Task and Synthetic Profile Generation}
We defined the retrieval objective: given a corrupted query profile snippet, retrieve the canonical record from a base index of 1,000,000 clean B2B profiles. 
We synthetically generate profiles consisting of a first name, last name, organization, work email, personal email, and country.
We evaluate models across six distinct query corruption buckets:
\begin{itemize}
    \item \textbf{Pristine}: No corruption.
    \item \textbf{Missing Firstname}: Omit the first name.
    \item \textbf{Missing Email/Company}: Omit the email address and company completely.
    \item \textbf{Typo Name}: Levenshtein distance mutations introduced on names.
    \item \textbf{Domain Mismatch}: Swap a generic `.gmail` domain for a corporate one.
    \item \textbf{Swapped Attributes}: Swap first name and last name fields.
\end{itemize}

\subsection{Models and Serialization}
We evaluate the BM25 lexical baseline alongside four dense retriever architectures: all-MiniLM-L6-v2 (22M parameters) \cite{reimers2019sentence}, bge-small-en-v1.5 (33M) \cite{wang-etal-2022-text}, nomic-embed-text-v1.5 (137M), and gte-modernbert-base (149M).
Records are serialized into a minimal positional structure (the "pipe" format):
\texttt{Jonathan Smith | Google Inc | jonathan.smith@google.com | USA}

\subsection{Fine-Tuning Procedure}
We construct a triplets dataset consisting of an anchor (a corrupted query), a positive (the canonical complete record), and a negative (either a hard negative sharing a domain or last name, or a random negative). Hard negative ratios are escalated across epochs (curriculum learning).
The models are trained using Multiple Negatives Ranking Loss (MNRL) \cite{henderson2017efficient} wrapped in MatryoshkaLoss, evaluating dimensionalities of [768, 512, 256, 128, 64]. 

\section{Experiments}

\subsection{Latency and Quantization Ablation}
We derive heavily compressed candidate indices directly from the core FP32 indexes to validate dimensional and quantization effects. We test FP32, INT8, and Binary quantization across varying dimensionalities without re-running the heavy neural encoding pipeline. Dense vector recall is scored using exact-nearest neighbor Search for accurate measurements before introducing ANN degradation. Latencies reported are purely system retrieve (LanceDB p50 limits).

\subsection{Results: Missing Information and Typos}

\begin{table}[h]
\centering
\caption{Recall@10 Across Model Configurations and Corruption Buckets}
\begin{tabular}{lcccccc}
\toprule
\textbf{Model Config} & \textbf{Pristine} & \textbf{Missing} ($fn$) & \textbf{Missing} ($em/co$) & \textbf{Typo} & \textbf{Domain}  & \textbf{Overall R@10} \\ \midrule
BM25 (Lexical Base) & \textbf{1.000} & \textbf{1.000} & 0.750 & \textbf{1.000} & \textbf{1.000} & 0.958 \\ 
GTE-ModernBERT Base & 1.000 & 0.919 & 0.747 & 0.994 & 0.990 & 0.941 \\ 
\textbf{GTE-ModernBERT FT} & 1.000 & 0.997 & \textbf{0.798} & 1.000 & 1.000 & \textbf{0.966} ($+$0.7pp) \\ 
BGE-Small FT (256-int8) & 1.000 & 0.984 & 0.672 & 0.986 & 1.000 & 0.942 ($-$1.6pp) \\
MiniLM-L6 FT (128-int8) & 1.000 & 0.979 & 0.647 & 0.979 & 1.000 & 0.934 ($-$2.4pp) \\
Nomic-v1.5 Base & 0.997 & 0.930 & 0.488 & 0.966 & 0.916 & 0.882 ($-$7.6pp) \\
Nomic-v1.5 FT & 0.979 & 0.959 & 0.152 & 0.948 & 0.895 & 0.817 ($-$14.2pp) \\
\bottomrule
\end{tabular}
\end{table}

\textbf{Winning Heavily Corrupted Queries.}
GTE-ModernBERT (149M parameters) fine-tuned on the task achieves a positive net-gain over BM25 in overall Recall@10 (0.966 vs 0.958). The fundamental advantage emerges in the most aggressive corruption condition (omitting email and company), where the GTE model leaps $+4.8$pp over the lexical baseline. 

\textbf{Catastrophic Forgetting of Implicit Prefixes.}
Fine-tuning the 137M parameter \texttt{nomic-v1.5} produced catastrophic behavioral drift. Nomic natively requires \texttt{"search\_query: "} strings appended. Applying generic MNRL optimization strictly mapped the target domain string positions while destroying the universal routing mechanism derived from its original causal contrastive objective. Recall on missing fields plummeted from 48.8\% (zero-shot) to 15.2\%.

\textbf{Index Volume Optimization.}
While GTE-ModernBERT produces the strongest R@10 metrics, it triggers a baseline 24.12ms p50 latency vs BM25 at 3.25ms. Small-footprint alternatives scaled well. We observe that \texttt{minilm\_l6} retains 0.934 R@10 under aggressive 128-dimensional INT8 quantization. This condenses the representation payload heavily, yielding sub-7ms retrieval latency.

\section{Conclusion}

We demonstrate that for massive, structurally erratic data formats, dense retriever models offer substantive recall improvement over exact lexical matches explicitly on records heavily degraded by missing organizational contexts. Fine-tuning an explicit base framework such as ModernBERT offers optimal results, whereas rigid instruction-tuned architectures collapse when exposed entirely to generic structural embeddings objectives. We offer a methodology to strictly optimize index sizing through Matryoshka learning and integer derivation pipelines to match required real-time latencies safely. 

\bibliographystyle{plain}
\bibliography{refs.bib}

\end{document}
